% eksperymentalne tłumaczenie na włoski

%

\section{Uszkodzone przyrządy} % lang-pl
\section{Strumenti danneggiati} % lang-it

Problemy od \ref{broken_ruler_compass_hartshorne_start} do \ref{broken_ruler_compass_hartshorne_end} pochodzą z książki Hartshorne'a \cite[s. 25, 26]{hartshorne2000}. % lang-pl
I problemi da \ref{broken_ruler_compass_hartshorne_start} a \ref{broken_ruler_compass_hartshorne_end} provengono dal libro di Hartshorne \cite[p. 25, 26]{hartshorne2000}. % lang-it

\begin{geoconstruction}
[połamana linijka] % lang-pl
[riga spezzata] % lang-it
\label{broken_ruler_compass_hartshorne_start}%
\index{linijka!połamana}%
    Dane są dwa punkty $A$ i $B$ na płaszczyźnie, odległe od siebie o około trzy nible. % lang-pl
    Mając do dyspozycji fragment linijki o długości jednej nibli oraz sprawny cyrkiel, narysować odcinek $AB$. % lang-pl
    Sono dati due punti $A$ e $B$ nel piano, distanti circa tre nibli. % lang-it
    Disponendo di un frammento di riga lungo una nible e di un compasso funzionante, tracciare il segmento $AB$. % lang-it
\end{geoconstruction}
% Hartshorne s. 25
O problemie pisze Eves \cite[s. 181]{eves1_1972}. % lang-pl
Del problema scrive Eves \cite[p. 181]{eves1_1972}. % lang-it

\begin{geoconstruction}
[zardzewiały cyrkiel] % lang-pl
[compasso arrugginito] % lang-it
    \index{cyrkiel!zardzewiały}
    Dane są dwa punkty $A$ i $B$ na płaszczyźnie, odległe od siebie o około pięć nibli. % lang-pl
    Mając do dyspozycji zardzewiały cyrkiel, którym można kreślić jedynie okręgi o promieniu dwóch nibli, skonstruować trójkąt równoboczny oparty o bok $AB$. % lang-pl
    Sono dati due punti $A$ e $B$ nel piano, distanti circa cinque nibli. % lang-it
    Disponendo di un compasso arrugginito con cui è possibile tracciare soltanto cerchi di raggio pari a due nibli, costruire un triangolo equilatero avente $AB$ come lato. % lang-it
\end{geoconstruction}

O zardzewiałym cyrklu będzie rozmyślać Abu al-Wafa Buzjani, matematyk perski, % lang-pl
Sul compasso arrugginito rifletterà Abu al-Wafa Buzjani, matematico persiano, % lang-it
\index[persons]{Buzjani, Abu al-Wafa}%
a pod koniec piętnastego wieku zastosują go w konstrukcjach Leonardo da Vinci oraz Albrecht Dürer. % lang-pl
e alla fine del XV secolo esso verrà impiegato nelle costruzioni di Leonardo da Vinci e Albrecht Dürer. % lang-it
\index[persons]{Dürer, Albrech}%
\index[persons]{da Vinci, Leonardo}%
% ŹRÓDŁO: https://en.wikipedia.org/wiki/Poncelet-Steiner_theorem#History

\begin{geoconstruction}
    Dany jest punkt $A$ leżący na prostej $l$. % lang-pl
    Skonstruować prostą prostopadłą do $l$ przechodzącą przez $A$ przy użyciu linijki i zardzewiałego cyrkla. % lang-pl
    È dato un punto $A$ appartenente alla retta $l$. % lang-it
    Costruire la retta perpendicolare a $l$ passante per $A$ usando una riga e un compasso arrugginito. % lang-it
\end{geoconstruction}
% Hartshorne s. 25

\begin{geoconstruction}
    Dany jest punkt $A$ leżący ponad cztery nible od prostej $l$. % lang-pl
    Skonstruować prostą prostopadłą do $l$ przechodzącą przez $A$ przy użyciu linijki i zardzewiałego cyrkla. % lang-pl
    È dato un punto $A$ a più di quattro nibli dalla retta $l$. % lang-it
    Costruire la retta perpendicolare a $l$ passante per $A$ usando una riga e un compasso arrugginito. % lang-it
\end{geoconstruction}
% Hartshorne s. 25

\begin{geoconstruction}
    Dane są trzy niewspółliniowe punkty $A$, $B$ oraz $C$. % lang-pl
    Skonstrować punkt $D$ na prostej $AC$ tak, żeby odcinki $AD$ oraz $AB$ były równej długości, przy użyciu linijki i zardzewiałego cyrkla. % lang-pl
    Sono dati tre punti non allineati $A$, $B$ e $C$. % lang-it
    Costruire un punto $D$ sulla retta $AC$ in modo che i segmenti $AD$ e $AB$ abbiano la stessa lunghezza, usando una riga e un compasso arrugginito. % lang-it
\end{geoconstruction}
% Hartshorne s. 26

\begin{geoconstruction}
    Dany jest odcinek $AB$ o długości ponad dwóch nibli oraz prosta $l$, która nie przechodzi przez końce odcinka. % lang-pl
    Skonstrować punkt $C$ na prostej $l$ tak, żeby odcinki $AB$ oraz $AC$ były równej długości, przy użyciu linijki i zardzewiałego cyrkla. % lang-pl
    Sono dati un segmento $AB$ di lunghezza superiore a due nibli e una retta $l$ che non passa per gli estremi del segmento. % lang-it
    Costruire un punto $C$ sulla retta $l$ in modo che i segmenti $AB$ e $AC$ abbiano la stessa lunghezza, usando una riga e un compasso arrugginito. % lang-it
\end{geoconstruction}
% Hartshorne s. 26

\begin{proposition}
\label{broken_ruler_compass_hartshorne_end}%
    Każdą konstrukcję geometryczną, którą można wykonać cyrklem i linijką, można powtórzyć przy użyciu zardzewiałego cyrkla i linijki. % lang-pl
    Qualsiasi costruzione geometrica realizzabile con riga e compasso può essere ripetuta usando una riga e un compasso arrugginito. % lang-it
\end{proposition}
% Hartshorne s. 26

Wynik ten zostanie odkryty relatywnie wcześnie. % lang-pl
Pierwszy pokaże go Lodovico Ferrari, uczeń Gerolamo Cardano i odkrywca metody rozwiązywania równań czwartego stopnia w pojedynku przeciwko Niccolò Fontanie Tartaglii. % lang-pl
Questo risultato sarà scoperto relativamente presto. % lang-it
Il primo a dimostrarlo sarà Lodovico Ferrari, allievo di Gerolamo Cardano e scopritore del metodo per risolvere le equazioni di quarto grado nella celebre disputa con Niccolò Fontana Tartaglia. % lang-it
\index[persons]{Ferrari, Lodovico}%
\index[persons]{Cardano, Gerolamo}%
W przeciągu dziesięciu lat wyczyn ten powtórzą Cardano, Tartaglia oraz uczeń Tartaglii, Benedetti. % lang-pl
Nel giro di dieci anni il risultato sarà ottenuto nuovamente da Cardano, Tartaglia e dall'allievo di Tartaglia, Benedetti. % lang-it
\index[persons]{Tartaglia, Nicolo}%
\index[persons]{Benedetti, Giambattista}%

Ciąg dalszy to oczywiście twierdzenie \ref{thm:poncelet_steiner}, że zardzewiałego cyrkla wystarczy użyć raz. % lang-pl
Il seguito naturale è naturalmente il teorema \ref{thm:poncelet_steiner}, secondo cui è sufficiente usare il compasso arrugginito una sola volta. % lang-it

%